\chapter{Deformations and Moduli}\label{deformations}
\chapterstatus{draft}

\section{Introduction}\label{deformations:section-introduction}

A compact Riemann surface of genus $1$ is a torus $\CC/\Lambda$, and different lattices can give non-isomorphic Riemann surfaces, although all of them are diffeomorphic to $S^1 \times S^1$. So a fixed
smooth manifold can carry many complex structures, and they vary continuously. Riemann already counted that the compact Riemann surfaces of genus $g \geq 2$ depend on $3g - 3$ complex parameters, his
``moduli''. Two questions arise: how do complex structures vary in families (\emph{deformation theory}), and is there a space whose points are the isomorphism classes, a \emph{moduli space}? This chapter
gives a first introduction to both.
\begin{enumerate}
\item Families of compact complex manifolds, and Ehresmann's theorem: in a family over a connected base all fibres are diffeomorphic
(Section~\ref{deformations:section-families}).
\item The moduli space of elliptic curves: isomorphism classes of complex tori of dimension one correspond to points of $\mathbb{H}/\SL_2(\ZZ)$, which has an explicit fundamental domain
(Section~\ref{deformations:section-elliptic}).
\item The Kodaira--Spencer map, which sends a tangent vector to the base of a family to a class in $H^1(X, T_X)$, the space of \emph{first-order deformations}; we compute it for the family of elliptic curves and
show that $\PP^1$ is rigid (Section~\ref{deformations:section-kodaira-spencer}).
\item An overview of the main theorems: Kuranishi's versal family, obstructions in $H^2(X, T_X)$, the Bogomolov--Tian--Todorov theorem, stability of K\"ahler manifolds, and moduli spaces of curves
(Section~\ref{deformations:section-kuranishi}).
\end{enumerate}
Deformation theory is where Hodge theory becomes variational: the Hodge decomposition of the fibres of a family varies holomorphically, and the derivative of this variation is computed from the
Kodaira--Spencer class (Chapter~\ref{vhs}). References: \cite[Chapter 6]{huybrechts-complex}, \cite[Chapter 9]{voisin-hodge-1}.

\noindent\textbf{Prerequisites.} Chapter~\ref{smooth-manifolds} (Smooth Manifolds), Chapter~\ref{complex-analysis} (Complex Analysis), Chapter~\ref{homotopy} (Homotopy and the Fundamental Group),
Chapter~\ref{complex-manifolds} (Complex Manifolds), Chapter~\ref{holomorphic-bundles} (Holomorphic Vector Bundles) and Chapter~\ref{dolbeault} (Dolbeault Cohomology).

Throughout, $T_X = T^{1,0}X$ denotes the holomorphic tangent bundle of a complex manifold $X$, and $H^1(X, T_X)$ its first Dolbeault cohomology $H^{0,1}_{\bar\partial}(X, T_X)$
(Definition~\ref{dolbeault:definition-dolbeault}); its elements are represented by $\bar\partial$-closed $(0,1)$-forms with values in $T_X$, that is, by expressions
$\sum_{j,k}a_{jk}\,d\bar z^j \otimes \partial/\partial z^k$ in local coordinates.

\section{Families of complex manifolds}\label{deformations:section-families}

\begin{definition}\label{deformations:definition-family}
A \emph{family of compact complex manifolds} over a connected complex manifold $B$ is a holomorphic map $\pi \colon \mathcal{X} \to B$ of complex manifolds that is proper, surjective and a submersion
(its differential is surjective at every point). The \emph{fibres} $X_t = \pi^{-1}(t)$ are compact complex submanifolds of $\mathcal{X}$. The family is \emph{trivial} if there is a biholomorphism
$\mathcal{X} \cong B \times X$ over $B$, and \emph{locally trivial} if every point of $B$ has a neighbourhood over which it is trivial. A \emph{deformation} of a compact complex manifold $X$ is a family
$\pi \colon \mathcal{X} \to B$ with a point $0 \in B$ and an isomorphism $X \cong X_0$.
\end{definition}

The fibres are complex submanifolds by the holomorphic version of the implicit function theorem: near a point of $\mathcal{X}$ there are holomorphic coordinates $(z^1, \ldots, z^n, t^1, \ldots, t^m)$ in
which $\pi(z, t) = t$. Indeed, if $t^1, \ldots, t^m$ are coordinates on $B$ near $\pi(x)$, the functions $t^k \circ \pi$ have independent differentials at $x$, and can be completed to a system of
coordinates by $n$ further linear functions $z^j$ of some chart; the map $(z, t)$ has invertible complex linear derivative, so by the inverse function theorem
(Theorem~\ref{real-analysis:theorem-inverse-function}) it is a diffeomorphism near $x$, and its inverse is holomorphic, since the derivative of the inverse is the inverse of a complex linear map.
We call such coordinates \emph{adapted} to $\pi$.

\begin{example}\label{deformations:example-families}
\begin{enumerate}
\item \emph{The family of elliptic curves.} Let $\mathbb{H} = \{\tau \in \CC : \operatorname{Im}\tau > 0\}$. The group $\ZZ^2$ acts on $\CC \times \mathbb{H}$ by
$(m, n)\cdot(z, \tau) = (z + m + n\tau, \tau)$, freely and properly discontinuously, by holomorphic maps. The quotient $\mathcal{E} = (\CC \times \mathbb{H})/\ZZ^2$ is a complex manifold, the projection
to $\mathbb{H}$ is a family, and its fibre over $\tau$ is the complex torus $E_\tau = \CC/(\ZZ + \ZZ\tau)$. Every fibre is diffeomorphic to $\RR^2/\ZZ^2$, by $(x, y) \mapsto x + \tau y$.
\item \emph{Plane curves.} The homogeneous polynomials of degree $d$ in three variables form a vector space of dimension $\binom{d+2}{2}$; those defining a smooth curve form an open subset $U$ of it, and
$\{(F, [x]) \in U \times \PP^2 : F(x) = 0\} \to U$ is a family of compact Riemann surfaces of genus $\frac{(d-1)(d-2)}{2}$. Since $F$ and $\lambda F$, and $F$ and $F \circ A$ for $A \in \GL_3(\CC)$, give
isomorphic curves, one expects about $\binom{d+2}{2} - 9$ effective parameters, which is $1$ for cubics and $6 = 3 \cdot 3 - 3$ for quartics, in accordance with Riemann's count for genus $1$ and $3$.
\item \emph{Hopf surfaces.} For $\alpha \in \CC$ with $|\alpha| > 1$, the group generated by $z \mapsto \alpha z$ acts freely and properly discontinuously on $\CC^2 \setminus \{0\}$, and the quotient $X_\alpha$
is a compact complex surface. These surfaces form a family over $\{|\alpha| > 1\}$ (the quotient of $(\CC^2 \setminus \{0\}) \times \{|\alpha| > 1\}$ by $(z, \alpha) \mapsto (\alpha z, \alpha)$). Every
$X_\alpha$ is diffeomorphic to $S^1 \times S^3$ (Exercise~\ref{deformations:exercise-hopf}); as $b_1(S^1 \times S^3) = 1$ is odd, no Hopf surface is K\"ahler (Chapter~\ref{hodge-decomposition}).
\item \emph{A trivial and a non-trivial family with isomorphic fibres.} The projection $B \times X \to B$ is trivial. On the other hand, let $E = E_i$ and let $\ZZ$ act on $\CC \times E$ by
$k\cdot(t, z) = (t + k, i^kz)$; the quotient is a family over $\CC/\ZZ \cong \CC^*$ all of whose fibres are isomorphic to $E$. It is not trivial: going once around the base induces the automorphism
$z \mapsto iz$ of $E$, which acts on $H_1(E; \ZZ) = \ZZ[i]$ as multiplication by $i$, whereas in a product the transport along a loop acts trivially on homology. Such families are \emph{isotrivial}. They are locally trivial; in contrast, Remark~\ref{deformations:remark-jumping} describes a family whose fibres are all isomorphic except one.
\end{enumerate}
\end{example}

The fibres of a family over a connected base are all diffeomorphic, so deformation theory studies the complex structures on a fixed smooth manifold. This is a theorem about smooth maps.

\begin{theorem}[Ehresmann]\label{deformations:theorem-ehresmann}
A proper surjective submersion $f \colon M \to B$ of smooth manifolds is a locally trivial fibre bundle: every $b \in B$ has a neighbourhood $U$ with a diffeomorphism
$f^{-1}(U) \cong U \times f^{-1}(b)$ over $U$.
\end{theorem}

\begin{proof}
Take $U$ a coordinate ball centred at $b$ with coordinates $x^1, \ldots, x^m$. Since $df$ is surjective, the coordinate vector fields $\partial/\partial x^i$ lift locally to vector fields on
$f^{-1}(U)$, and a partition of unity (Theorem~\ref{smooth-manifolds:theorem-partitions-of-unity}) glues local lifts to vector fields $v_i$ on $f^{-1}(U)$ with $df(v_i) = \partial/\partial x^i$,
since the condition is affine. As $f$ is proper, the flows of the $v_i$ exist as long as the flows of $\partial/\partial x^i$ stay in $U$
(Theorem~\ref{smooth-manifolds:theorem-flow}: an integral curve cannot leave a compact set in finite time, and $f^{-1}$ of a compact subset of $U$ is compact). The map
$U \times f^{-1}(b) \to f^{-1}(U)$, $(x, y) \mapsto \phi^1_{x^1} \circ \cdots \circ \phi^m_{x^m}(y)$, composing the flows for the times $x^i$, is then a diffeomorphism over $U$, with inverse obtained by
flowing back.
\end{proof}

\begin{corollary}\label{deformations:corollary-fibres-diffeomorphic}
All fibres of a family of compact complex manifolds over a connected base are diffeomorphic. In particular their Betti numbers agree.
\end{corollary}

\begin{proof}
By Theorem~\ref{deformations:theorem-ehresmann}, the set of $t \in B$ whose fibre is diffeomorphic to a given fibre $X_{t_0}$ is open, and so is its complement. As $B$ is connected, it is all of $B$.
\end{proof}

The trivialisation is only smooth: Ehresmann's theorem says nothing about the complex structures, which genuinely vary, as the next section shows for elliptic curves.

\section{The moduli space of elliptic curves}\label{deformations:section-elliptic}

A \emph{lattice} in $\CC$ is a subgroup $\Lambda = \ZZ\omega_1 + \ZZ\omega_2$ generated by two $\RR$-linearly independent numbers; the quotient $\CC/\Lambda$ is a compact Riemann surface, an \emph{elliptic curve},
diffeomorphic to $S^1 \times S^1$. Every compact Riemann surface of genus $1$ is of this form, by the uniformisation theorem (Theorem~\ref{complex-manifolds:theorem-uniformisation}), but we shall not need this.
We determine when two such tori are isomorphic.

\begin{proposition}\label{deformations:proposition-tori-isomorphic}
Let $\Lambda$, $\Lambda^{\prime} \subseteq \CC$ be lattices. Every holomorphic map $f \colon \CC/\Lambda \to \CC/\Lambda^{\prime}$ is induced by an affine map $z \mapsto \alpha z + \beta$ with $\alpha\Lambda \subseteq \Lambda^{\prime}$.
Consequently $\CC/\Lambda \cong \CC/\Lambda^{\prime}$ if and only if $\Lambda^{\prime} = \alpha\Lambda$ for some $\alpha \in \CC^*$.
\end{proposition}

\begin{proof}
The projection $\CC \to \CC/\Lambda^{\prime}$ is a covering map and $\CC$ is simply connected, so $f \circ p \colon \CC \to \CC/\Lambda^{\prime}$, with $p \colon \CC \to \CC/\Lambda$, lifts to a continuous map
$F \colon \CC \to \CC$ (Proposition~\ref{homotopy:proposition-lifting-criterion}), which is holomorphic because the projections are local biholomorphisms. For $\lambda \in \Lambda$, $F(z + \lambda) - F(z)$ lies in
the discrete set $\Lambda^{\prime}$ and depends continuously on $z$, so it is a constant $c(\lambda)$. Hence $F^{\prime}(z + \lambda) = F^{\prime}(z)$: the entire function $F^{\prime}$ is $\Lambda$-periodic, hence bounded
(its values are those on a compact fundamental parallelogram), hence constant by Liouville's theorem (Corollary~\ref{complex-analysis:corollary-liouville}). So $F(z) = \alpha z + \beta$, and
$\alpha\lambda = c(\lambda) \in \Lambda^{\prime}$. If $f$ is an isomorphism, the same applied to $f^{-1}$ gives $\alpha^{-1}\Lambda^{\prime} \subseteq \Lambda$, so $\alpha\Lambda = \Lambda^{\prime}$ with $\alpha \neq 0$. Conversely,
$z \mapsto \alpha z$ induces an isomorphism $\CC/\Lambda \to \CC/\alpha\Lambda$.
\end{proof}

Every lattice has a basis $\omega_1, \omega_2$ with $\operatorname{Im}(\omega_2/\omega_1) > 0$ (exchange the two vectors if necessary), and then $\Lambda = \omega_1(\ZZ + \ZZ\tau)$ with $\tau = \omega_2/\omega_1 \in \mathbb{H}$.
So every elliptic curve is isomorphic to some $E_\tau = \CC/(\ZZ + \ZZ\tau)$. The group $\SL_2(\ZZ)$ acts on $\mathbb{H}$ by Möbius transformations,
\[
\begin{pmatrix}a & b\\ c & d\end{pmatrix}\cdot\tau = \frac{a\tau + b}{c\tau + d}, \qquad \operatorname{Im}\frac{a\tau + b}{c\tau + d} = \frac{\operatorname{Im}\tau}{|c\tau + d|^2},
\]
the formula for the imaginary part following from $(a\tau + b)(c\bar\tau + d) = ac|\tau|^2 + ad\tau + bc\bar\tau + bd$ and $ad - bc = 1$.

\begin{theorem}\label{deformations:theorem-elliptic-moduli}
$E_\tau \cong E_{\tau^{\prime}}$ if and only if $\tau^{\prime} = \gamma\cdot\tau$ for some $\gamma \in \SL_2(\ZZ)$. So the isomorphism classes of elliptic curves correspond bijectively to the points of
$\mathbb{H}/\SL_2(\ZZ)$. Every orbit meets the set
\[
\mathcal{F} = \{\tau \in \mathbb{H} : |\operatorname{Re}\tau| \leq \tfrac12, \ |\tau| \geq 1\},
\]
and two distinct points of $\mathcal{F}$ are in the same orbit only if they lie on the boundary and are exchanged by $\tau \mapsto \tau + 1$ or by $\tau \mapsto -1/\tau$.
\end{theorem}

\begin{proof}
By Proposition~\ref{deformations:proposition-tori-isomorphic}, $E_\tau \cong E_{\tau^{\prime}}$ if and only if $\ZZ + \ZZ\tau^{\prime} = \alpha(\ZZ + \ZZ\tau)$ for some $\alpha$. Then $\alpha\tau$ and $\alpha$ form a basis of
$\ZZ + \ZZ\tau^{\prime}$, as do $\tau^{\prime}$ and $1$; so $\tau^{\prime} = \alpha(a\tau + b)$ and $1 = \alpha(c\tau + d)$ with $\bigl(\begin{smallmatrix}a & b\\ c & d\end{smallmatrix}\bigr) \in \GL_2(\ZZ)$, and dividing,
$\tau^{\prime} = (a\tau + b)/(c\tau + d)$. Since both $\tau$ and $\tau^{\prime}$ lie in $\mathbb{H}$, the formula for the imaginary part, valid for $\det = \pm1$ with the factor $\det$, forces $ad - bc = 1$. Conversely, if
$\tau^{\prime} = \gamma\cdot\tau$, then $\alpha = (c\tau + d)^{-1}$ satisfies $\alpha(\ZZ + \ZZ\tau) = \ZZ + \ZZ\tau^{\prime}$, as $\gamma$ is invertible over $\ZZ$.

\emph{Every orbit meets $\mathcal{F}$.} Recall that $\SL_2(\ZZ)$ contains $T = \bigl(\begin{smallmatrix}1 & 1\\ 0 & 1\end{smallmatrix}\bigr)$, acting by $\tau \mapsto \tau + 1$, and
$S = \bigl(\begin{smallmatrix}0 & -1\\ 1 & 0\end{smallmatrix}\bigr)$, acting by $\tau \mapsto -1/\tau$. Let $\tau \in \mathbb{H}$. The numbers $|c\tau + d|$ for $(c, d) \in \ZZ^2 \setminus \{0\}$ are bounded below by a positive constant and take only finitely many values below
any bound (the lattice $\ZZ + \ZZ\tau$ is discrete), so $\operatorname{Im}(\gamma\cdot\tau) = \operatorname{Im}\tau/|c\tau + d|^2$ attains a maximum over $\gamma \in \SL_2(\ZZ)$; choose $\gamma$ attaining it and a
power $T^k \colon \tau \mapsto \tau + k$ with $|\operatorname{Re}(T^k\gamma\cdot\tau)| \leq \frac12$. The point $\sigma = T^k\gamma\cdot\tau$ still has maximal imaginary part. If $|\sigma| < 1$, then $S \colon \sigma \mapsto -1/\sigma$
would give $\operatorname{Im}(-1/\sigma) = \operatorname{Im}\sigma/|\sigma|^2 > \operatorname{Im}\sigma$, a contradiction. So $\sigma \in \mathcal{F}$.

\emph{Points of $\mathcal{F}$ in the same orbit.} Let $\tau = x + iy$ and $\gamma\cdot\tau$ both lie in $\mathcal{F}$, with $\gamma = \bigl(\begin{smallmatrix}a & b\\ c & d\end{smallmatrix}\bigr) \neq \pm1$.
Exchanging the roles of $\tau$ and $\gamma\cdot\tau$ (and replacing $\gamma$ by $\gamma^{-1}$) we may assume $\operatorname{Im}(\gamma\cdot\tau) \geq \operatorname{Im}\tau$, that is, $|c\tau + d| \leq 1$. Write
$\rho = e^{2\pi i/3} = -\frac12 + \frac{\sqrt3}{2}i$ and $\rho + 1 = -\bar\rho = e^{\pi i/3}$ for the two corners of $\mathcal{F}$. For $\tau \in \mathcal{F}$ we have $y^2 = |\tau|^2 - x^2 \geq 1 - \frac14$, so
$y \geq \sqrt3/2$, and $|c|\sqrt3/2 \leq |c|y = |\operatorname{Im}(c\tau + d)| \leq 1$ forces $c \in \{0, \pm1\}$.
\begin{itemize}
\item If $c = 0$, then $ad = 1$, so $a = d = \pm1$ and $\gamma\cdot\tau = \tau \pm b$ with $b \neq 0$. Both $\tau$ and $\tau \pm b$ have real part in $[-\frac12, \frac12]$, so $|b| = 1$ and $x = \pm\frac12$: the two points
lie on the vertical sides and are exchanged by $\tau \mapsto \tau \pm 1$.
\item If $c = \pm1$, replacing $\gamma$ by $-\gamma$ (which acts in the same way) we may take $c = 1$, so $|\tau + d| \leq 1$. Now $|\tau + d|^2 = (x + d)^2 + y^2$. If $|d| \geq 2$ this is at least
$(\frac32)^2 > 1$. If $d = \pm1$, it is $|\tau|^2 \pm 2x + 1 \geq 1 \pm 2x + 1 \geq 1$, with equality only if $|\tau| = 1$ and $x = \mp\frac12$, that is, $\tau = \rho$ for $d = 1$ and $\tau = \rho + 1$ for $d = -1$.
If $d = 0$, it is $|\tau|^2 \geq 1$, with equality only if $|\tau| = 1$.

\emph{Case $d = 0$, $|\tau| = 1$.} Then $-b = ad - bc = 1$ and $\gamma\cdot\tau = (a\tau - 1)/\tau = a - 1/\tau = a - \bar\tau$, of real part $a - x$. This lies in $[-\frac12, \frac12]$ only if $a = 0$, or $a = \pm1$ and
$x = \pm\frac12$. If $a = 0$, then $\gamma\cdot\tau = -1/\tau$, and $\tau$, $-1/\tau = -\bar\tau$ are points of the arc $|\tau| = 1$ exchanged by $\tau \mapsto -1/\tau$. If $a = \pm1$, then $x = a/2$, so
$\tau$ is a corner, and $\gamma\cdot\tau = a - \bar\tau = (a - x) + iy = \tau$.

\emph{Case $d = \pm1$, $\tau$ a corner.} Then $b = ad - 1$ and $\gamma\cdot\tau = a - (\tau + d)^{-1}$. For $\tau = \rho$, $d = 1$: $\rho + 1 = -\bar\rho$ and $(\rho + 1)^{-1} = -\rho$, so $\gamma\cdot\rho = a + \rho$,
which is in $\mathcal{F}$ only for $a \in \{0, 1\}$, giving $\rho$ itself or $\rho + 1$. For $\tau = \rho + 1$, $d = -1$, similarly $\gamma\cdot\tau = a - \rho^{-1} = a - \bar\rho = a + \rho + 1$, giving $\rho + 1$ or $\rho$.
\end{itemize}
In all cases either $\gamma\cdot\tau = \tau$ or the two points lie on the boundary of $\mathcal{F}$ and are exchanged by $\tau \mapsto \tau \pm 1$ or by $\tau \mapsto -1/\tau$.
\end{proof}

So the \emph{moduli space of elliptic curves} is $\mathbb{H}/\SL_2(\ZZ)$, which is $\mathcal{F}$ with its boundary glued: the vertical sides are identified by $\tau \mapsto \tau + 1$ and the two halves of the
arc $|\tau| = 1$ by $\tau \mapsto -1/\tau$. Topologically it is a plane. The case analysis in the proof also shows that the stabiliser of a point of $\mathcal{F}$ in $\SL_2(\ZZ)$ is $\{\pm1\}$, except at
$\tau = i$ (stabiliser of order $4$, generated by $S$) and at the corners $\rho$, $\rho + 1$ (order $6$; for $\rho$ it is generated by $ST$, as $ST\cdot\tau = -1/(\tau + 1)$); correspondingly every elliptic curve has the automorphism $z \mapsto -z$, $E_i$ has in addition $z \mapsto iz$, and $E_\rho$ has $z \mapsto -\rho z$, a sixth root of
unity. These two curves have \emph{complex multiplication} by $\ZZ[i]$ and $\ZZ[\rho]$, the Gaussian and Eisenstein integers.

\begin{remark}\label{deformations:remark-elliptic-moduli}
\begin{enumerate}
\item \emph{The $j$-invariant.} There is a holomorphic function $j \colon \mathbb{H} \to \CC$, invariant under $\SL_2(\ZZ)$, that induces a bijection $\mathbb{H}/\SL_2(\ZZ) \to \CC$; with the Eisenstein series
$g_2(\tau) = 60\sum{}^{\prime}(m + n\tau)^{-4}$ and $g_3(\tau) = 140\sum{}^{\prime}(m + n\tau)^{-6}$ it is $j = 1728g_2^3/(g_2^3 - 27g_3^2)$, with $j(i) = 1728$ and $j(\rho) = 0$ \cite[Chapter VII]{serre-arithmetic}. So the moduli
space of elliptic curves is the affine line, the \emph{$j$-line}.
\item \emph{Quadratic forms.} The reduction of $\tau$ into $\mathcal{F}$ is the reduction theory of positive definite binary quadratic forms (Theorem~\ref{number-theory:theorem-reduction}) in another guise:
the form $|x + y\tau|^2 = x^2 + 2\operatorname{Re}(\tau)xy + |\tau|^2y^2$ satisfies $|2\operatorname{Re}\tau| \leq 1 \leq |\tau|^2$, the reduction inequalities $|b| \leq a \leq c$, exactly when $\tau \in \mathcal{F}$. For a
positive definite form $ax^2 + bxy + cy^2$ of discriminant $D < 0$, the root $\tau = \frac{-b + i\sqrt{|D|}}{2a}$ of $a\tau^2 + b\tau + c$ in $\mathbb{H}$ gives an elliptic curve with complex multiplication, and
$j(\tau)$ is an algebraic integer of degree the class number $h(D)$; for class number one it is a rational integer, for instance $j\bigl(\frac{1 + i\sqrt{163}}{2}\bigr) = -640320^3$, which explains why
$e^{\pi\sqrt{163}}$ is so close to an integer (Remark~\ref{number-theory:remark-heegner}) \cite{cox-primes}.
\item \emph{Not a fine moduli space.} There is no family over $\mathbb{H}/\SL_2(\ZZ)$ whose fibre over each point is the corresponding curve and which induces all families: the isotrivial family of
Example~\ref{deformations:example-families}(4) has constant ``moduli point'' $[i]$ without being trivial. The automorphisms of the curves are the obstruction. Such a space, which classifies the objects but
does not carry a universal family, is a \emph{coarse moduli space}.
\end{enumerate}
\end{remark}

\section{The Kodaira--Spencer map}\label{deformations:section-kodaira-spencer}

Let $\pi \colon \mathcal{X} \to B$ be a family, $0 \in B$ and $X = X_0$. We want to measure, to first order, how the complex structure of $X_t$ changes as $t$ moves away from $0$ in a direction $v \in T_0B$. Locally
over $B$ nothing happens: the family looks like a product $U \times \Delta$ near each point of $X$. The information is in how these local products are glued.

Choose holomorphic coordinates $t = (t^1, \ldots, t^m)$ on $B$ centred at $0$. By compactness of $X$ and the existence of adapted coordinates, after shrinking $B$ to a small polydisc $\Delta$ there are open
sets $\mathcal{W}_i$ covering $\pi^{-1}(\Delta)$, $i = 1, \ldots, N$, with adapted coordinates $(z_i, t)$, $z_i = (z_i^1, \ldots, z_i^n)$, identifying $\mathcal{W}_i$ with $U_i \times \Delta$ for open
$U_i \subseteq \CC^n$; the sets $X \cap \mathcal{W}_i$ cover $X$. On overlaps the coordinates are related by holomorphic \emph{transition maps}
\[
z_j = f_{ij}(z_i, t), \qquad f_{ik}(z_i, t) = f_{jk}(f_{ij}(z_i, t), t).
\]
For a tangent vector $v = \sum_lv^l\partial/\partial t^l \in T_0B$ define, on $X \cap \mathcal{W}_i \cap \mathcal{W}_j$, the holomorphic vector field
\[
\theta_{ij}(v) = \sum_{k=1}^n\frac{\partial f_{ij}^k}{\partial v}(z_i, 0)\,\frac{\partial}{\partial z_j^k}, \qquad \frac{\partial}{\partial v} = \sum_lv^l\frac{\partial}{\partial t^l},
\]
the velocity, at $t = 0$, of the point with fixed coordinate $z_i$, as seen in the coordinates $z_j$.

\begin{lemma}\label{deformations:lemma-kodaira-spencer-cocycle}
The vector fields $\theta_{ij} = \theta_{ij}(v)$ satisfy $\theta_{ik} = \theta_{ij} + \theta_{jk}$ on triple overlaps; in particular $\theta_{ii} = 0$ and $\theta_{ji} = -\theta_{ij}$. If the adapted coordinates are
changed to $z_i^{\prime} = g_i(z_i, t)$, the new vector fields are $\theta_{ij}^{\prime} = \theta_{ij} + \eta_j - \eta_i$, where $\eta_i$ is the holomorphic vector field $\sum_k\frac{\partial g_i^k}{\partial v}(z_i, 0)\,\partial/\partial z_i^{\prime k}$
on $X \cap \mathcal{W}_i$.
\end{lemma}

\begin{proof}
Differentiate $f_{ik}(z_i, t) = f_{jk}(f_{ij}(z_i, t), t)$ with respect to $v$ at $t = 0$, by the chain rule:
\[
\frac{\partial f_{ik}}{\partial v} = D_zf_{jk}\cdot\frac{\partial f_{ij}}{\partial v} + \frac{\partial f_{jk}}{\partial v}.
\]
The first term is the vector $\theta_{ij}$, written in the coordinates $z_j$, transformed into the coordinates $z_k$ by the Jacobian matrix $D_zf_{jk}$ of the change of coordinates; that is, it is the same vector
field $\theta_{ij}$ expressed in the coordinates $z_k$. So $\theta_{ik} = \theta_{ij} + \theta_{jk}$. Taking $i = j = k$ gives $\theta_{ii} = 2\theta_{ii}$, and $k = i$ gives $\theta_{ji} = -\theta_{ij}$. For the change of
coordinates, the new transition maps are $f_{ij}^{\prime} = g_j \circ f_{ij} \circ g_i^{-1}$, where $g_i^{-1}$ inverts $z_i \mapsto g_i(z_i, t)$ for fixed $t$; the same differentiation, applied to
$f_{ij}^{\prime}(g_i(z_i, t), t) = g_j(f_{ij}(z_i, t), t)$, gives $\eta_i + \theta_{ij}^{\prime} = \theta_{ij} + \eta_j$ as vector fields.
\end{proof}

A family of holomorphic vector fields $(\theta_{ij})$ with $\theta_{ik} = \theta_{ij} + \theta_{jk}$ is a \emph{\v{C}ech $1$-cocycle} of the sheaf of holomorphic vector fields
(Definition~\ref{sheaf-cohomology:definition-cech}), and those of the form $\eta_j - \eta_i$ are the coboundaries. We convert it into a Dolbeault class, which does not depend on the cover.

\begin{lemma}\label{deformations:lemma-cech-dolbeault}
Let $E$ be a holomorphic vector bundle on a complex manifold $X$, $(U_i)$ an open cover, and $\theta_{ij}$ holomorphic sections of $E$ over $U_i \cap U_j$ with $\theta_{ik} = \theta_{ij} + \theta_{jk}$.
\begin{enumerate}
\item There are smooth sections $\sigma_i$ of $E$ over $U_i$ with $\sigma_j - \sigma_i = \theta_{ij}$ on $U_i \cap U_j$; for instance $\sigma_i = \sum_k\rho_k\theta_{ki}$ for a partition of unity $(\rho_k)$ subordinate
to the cover.
\item The forms $\bar\partial\sigma_i$ agree on overlaps, and define a $\bar\partial$-closed form $\delta(\theta) \in \mathcal{A}^{0,1}(X, E)$, whose class in $H^{0,1}_{\bar\partial}(X, E)$ does not depend on the
choice of the $\sigma_i$.
\item If $\theta_{ij} = \eta_j - \eta_i$ with $\eta_i$ holomorphic on $U_i$, then $[\delta(\theta)] = 0$. If $(V_a)$ is a refinement of the cover, with $V_a \subseteq U_{\lambda(a)}$, then the restricted cocycle
$\theta_{\lambda(a)\lambda(b)}$ gives the same class.
\end{enumerate}
\end{lemma}

\begin{proof}
(1) Each $\rho_k\theta_{ki}$, extended by $0$ outside $\operatorname{supp}\rho_k$, is a smooth section on $U_i$, and the sum is locally finite. Then
$\sigma_j - \sigma_i = \sum_k\rho_k(\theta_{kj} - \theta_{ki}) = \sum_k\rho_k\theta_{ij} = \theta_{ij}$, using $\theta_{kj} = \theta_{ki} + \theta_{ij}$.

(2) On $U_i \cap U_j$, $\bar\partial\sigma_j - \bar\partial\sigma_i = \bar\partial\theta_{ij} = 0$ since $\theta_{ij}$ is holomorphic. The glued form is locally $\bar\partial$-exact, hence $\bar\partial$-closed. If $\sigma_i^{\prime}$ is another
choice, then $\sigma_i^{\prime} - \sigma_i = \sigma_j^{\prime} - \sigma_j$ on overlaps, so these differences glue to a global smooth section $s$, and the two forms differ by $\bar\partial s$.

(3) If $\theta_{ij} = \eta_j - \eta_i$, we may take $\sigma_i = \eta_i$, which are holomorphic, so $\delta(\theta) = 0$. For a refinement, the sections $\sigma_{\lambda(a)}|_{V_a}$ satisfy the required relation and
have the same $\bar\partial$.
\end{proof}

\begin{definition}\label{deformations:definition-kodaira-spencer}
The \emph{Kodaira--Spencer map} of the family $\pi$ at $0$ is
\[
\kappa \colon T_0B \to H^1(X, T_X), \qquad \kappa(v) = [\delta(\theta(v))].
\]
By Lemmas~\ref{deformations:lemma-kodaira-spencer-cocycle} and~\ref{deformations:lemma-cech-dolbeault} it does not depend on the choice of adapted coordinates (two systems of adapted coordinates
can be compared on a common refinement), and it is clearly $\CC$-linear in $v$. The space $H^1(X, T_X)$ is the space of \emph{first-order deformations} of $X$.
\end{definition}

\begin{proposition}\label{deformations:proposition-kodaira-spencer-properties}
\begin{enumerate}
\item For a trivial family $B \times X \to B$, $\kappa = 0$.
\item (Base change) If $\varphi \colon B^{\prime} \to B$ is holomorphic with $\varphi(0^{\prime}) = 0$, the pulled-back family $\mathcal{X}^{\prime} = B^{\prime} \times_B \mathcal{X} \to B^{\prime}$ has Kodaira--Spencer map
$\kappa^{\prime} = \kappa \circ d\varphi_{0^{\prime}}$.
\end{enumerate}
\end{proposition}

\begin{proof}
(1) With adapted coordinates $(z_i, t)$ coming from charts $z_i$ of $X$, the transition maps $f_{ij}(z_i, t) = f_{ij}(z_i)$ do not depend on $t$, so $\theta = 0$. (2) The fibre product
$\mathcal{X}^{\prime} = \{(s, x) \in B^{\prime} \times \mathcal{X} : \varphi(s) = \pi(x)\}$ is, in adapted coordinates $(z_i, t)$ of $\mathcal{X}$, the set of points $(s, z_i, t)$ with $t = \varphi(s)$; so it is a complex
submanifold of $B^{\prime} \times \mathcal{X}$ with coordinates $(z_i, s)$, adapted to the projection to $B^{\prime}$, and its transition maps are $z_j = f_{ij}(z_i, \varphi(s))$. Its fibre over $0^{\prime}$ is $X_0$. By the
chain rule, $\partial f_{ij}(z_i, \varphi(s))/\partial w = \partial f_{ij}/\partial(d\varphi_{0^{\prime}}(w))$ at $s = 0^{\prime}$, so $\theta^{\prime}_{ij}(w) = \theta_{ij}(d\varphi_{0^{\prime}}(w))$.
\end{proof}

So if $\kappa(v) \neq 0$, the family is not trivial near $0$, and more: its restriction to any holomorphic disc $\varphi \colon \Delta \to B$ through $0$ with tangent vector $v$ is not trivial, since by (2) its
Kodaira--Spencer map sends $\partial/\partial s$ to $\kappa(v) \neq 0$, while by (1) a trivial family has $\kappa = 0$.

\begin{example}\label{deformations:example-kodaira-spencer-elliptic}
\emph{The family of elliptic curves.} Let $\mathcal{E} \to \mathbb{H}$ be the family of Example~\ref{deformations:example-families}(1), and $\tau_0 \in \mathbb{H}$. Adapted coordinates are given by $(z, \tau)$ on small open
sets of $\CC \times \mathbb{H}$ that map injectively to $\mathcal{E}$. Two such charts differ by an element of $\ZZ^2$, so the transition maps are $z_j = z_i + m_{ij} + n_{ij}\tau$ with integers $m_{ij}$, $n_{ij}$,
and
\[
\theta_{ij}\Bigl(\frac{\partial}{\partial\tau}\Bigr) = n_{ij}\frac{\partial}{\partial z}.
\]
Here $\partial/\partial z$ is a global holomorphic vector field on $E = E_{\tau_0}$, invariant under translations. To find $\sigma_i$ as in Lemma~\ref{deformations:lemma-cech-dolbeault}, write $z = x + \tau_0y$ with $x, y$ real;
the function $y = \operatorname{Im}z/\operatorname{Im}\tau_0$ satisfies $y(z + m + n\tau_0) = y(z) + n$. So $\sigma_i = y(z_i)\,\partial/\partial z$ satisfies $\sigma_j - \sigma_i = n_{ij}\,\partial/\partial z = \theta_{ij}$, and
\[
\kappa\Bigl(\frac{\partial}{\partial\tau}\Bigr) = \Bigl[\frac{\partial y}{\partial\bar z}\,d\bar z \otimes \frac{\partial}{\partial z}\Bigr] = \frac{i}{2\operatorname{Im}\tau_0}\Bigl[d\bar z \otimes \frac{\partial}{\partial z}\Bigr],
\]
since $\operatorname{Im}z = (z - \bar z)/2i$ has $\partial/\partial\bar z$ equal to $-1/2i = i/2$. This class is nonzero. As $T_E$ is trivialised by $\partial/\partial z$, it suffices to see that $d\bar z$ is not
$\bar\partial g$ for a smooth function $g$ on $E$. In the real coordinates $(x, y)$ of $\RR^2/\ZZ^2$, $\partial/\partial\bar z = \alpha\,\partial/\partial x + \beta\,\partial/\partial y$ with constants $\alpha$, $\beta$ (as $x$ and $y$
are real linear functions of $z$, $\bar z$), so if $\partial g/\partial\bar z = 1$, integrating over $[0, 1]^2$ and using periodicity, $1 = \int_{[0,1]^2}(\alpha g_x + \beta g_y)\,dx\,dy = 0$ by Fubini's theorem and the
fundamental theorem of calculus, a contradiction. So $\kappa$ is injective at every point: the family of elliptic curves is nowhere locally trivial. Here $H^1(E, T_E)$ is one-dimensional
(Proposition~\ref{deformations:proposition-curves-deformations}), so $\kappa$ is an isomorphism.
\end{example}

The dimension of $H^1(X, T_X)$ can often be computed with Serre duality. For a compact Riemann surface $X$, $T_X^\vee = K_X$, and Serre duality (Theorem~\ref{dolbeault:theorem-serre-duality}) gives
\[
H^1(X, T_X) \cong H^0(X, K_X \otimes K_X)^\vee,
\]
the dual of the space of holomorphic \emph{quadratic differentials} $q = f(z)\,dz^2$.

\begin{proposition}\label{deformations:proposition-curves-deformations}
\begin{enumerate}
\item $H^1(\PP^1, T_{\PP^1}) = 0$: the Riemann sphere has no first-order deformations.
\item For an elliptic curve $E$, $H^1(E, T_E)$ is one-dimensional.
\end{enumerate}
\end{proposition}

\begin{proof}
(1) A holomorphic quadratic differential on $\PP^1$ is $f(z)\,dz^2$ with $f$ entire on the chart $\CC$; in the coordinate $w = 1/z$ near $\infty$, $dz = -w^{-2}dw$, so it is $f(1/w)w^{-4}\,dw^2$, which is holomorphic at
$w = 0$ only if $f(z)z^4$ stays bounded as $z \to \infty$. Then $f$ is bounded, hence constant by Liouville's theorem, and the constant is $0$. So $H^0(K^2) = 0$, and $H^1(T) = 0$ by Serre duality.

(2) $K_E$ is trivialised by the global form $dz$, so $H^0(E, K_E^2) \cong H^0(E, \mathcal{O}_E)$, the holomorphic functions on the compact connected $E$. Such a function $h$ attains its maximum modulus at some point, so it is constant near that point by the maximum modulus principle
(Theorem~\ref{complex-analysis:theorem-maximum-modulus}), and the set where $h$ equals this constant is open and closed, hence all of $E$. So $H^0(E, K_E^2)$ is one-dimensional.
\end{proof}

For a curve of genus $g \geq 2$, the Riemann--Roch theorem gives $\dim H^0(X, K_X^2) = 3g - 3$ (Chapter~\ref{curves}), so $\dim H^1(X, T_X) = 3g - 3$, Riemann's count of moduli. For a complex torus $\CC^n/\Lambda$,
$T_X$ is trivial and $\dim H^1(X, T_X) = n\dim H^{0,1}_{\bar\partial}(X) = n^2$; this matches the count of linear complex structures on $\RR^{2n}$: $\GL_{2n}(\RR)/\GL_n(\CC)$ has real dimension $4n^2 - 2n^2 = 2n^2$.

\begin{remark}\label{deformations:remark-maurer-cartan}
\emph{Deformations as a differential equation.} By Ehresmann's theorem a deformation of $X$ over a disc can be viewed as a family of complex structures on the fixed manifold $X$. A complex structure $J_t$
close to $J_0$ is determined by its space of $(0,1)$-vectors, which is the graph of a map $T^{0,1} \to T^{1,0}$, that is, by a form $\varphi(t) \in \mathcal{A}^{0,1}(X, T_X)$ with $\varphi(0) = 0$. The new
$\bar\partial$-operator is $\bar\partial + \varphi$ acting through the contraction of $\varphi$ with $\partial f$, and the integrability of $J_t$ (Theorem~\ref{complex-manifolds:theorem-newlander-nirenberg}) becomes
the \emph{Maurer--Cartan equation}
\[
\bar\partial\varphi + \tfrac12[\varphi, \varphi] = 0 .
\]
Writing $\varphi(t) = t\varphi_1 + t^2\varphi_2 + \cdots$, the first-order part satisfies $\bar\partial\varphi_1 = 0$, and $[\varphi_1]$ is the Kodaira--Spencer class of $\partial/\partial t$. The second-order part requires
$\bar\partial\varphi_2 = -\frac12[\varphi_1, \varphi_1]$: the class of $[\varphi_1, \varphi_1]$ in $H^2(X, T_X)$ is the first \emph{obstruction} to extending $\varphi_1$ to an actual deformation
\cite[Chapter 6]{huybrechts-complex}.
\end{remark}

\section{Kuranishi theory and moduli spaces}\label{deformations:section-kuranishi}

We state the main theorems of the theory; their proofs combine the Maurer--Cartan equation with the harmonic theory of the $\bar\partial$-Laplacian and power series estimates.

\begin{theorem}[Kodaira--Spencer, Kuranishi]\label{deformations:theorem-kuranishi}
Let $X$ be a compact complex manifold.
\begin{enumerate}
\item There is a deformation $\mathcal{X} \to (K, 0)$ of $X$ over a complex analytic space $K$, the \emph{Kuranishi space}, which is \emph{complete}: every deformation of $X$ over a germ $(B, 0)$ is obtained from it by
pull-back along a holomorphic map $(B, 0) \to (K, 0)$, whose differential at $0$ is determined by the Kodaira--Spencer map. The Zariski tangent space of $K$ at $0$ is $H^1(X, T_X)$.
\item $K$ is the zero set, near $0$, of a holomorphic map $H^1(X, T_X) \to H^2(X, T_X)$ vanishing to second order at $0$. In particular, if $H^2(X, T_X) = 0$, then $K$ is a neighbourhood of $0$ in
$H^1(X, T_X)$, and $X$ is \emph{unobstructed}.
\item If $H^0(X, T_X) = 0$, that is, $X$ has no holomorphic vector fields, the map in (1) is unique, and the family is \emph{universal}.
\item If $H^1(X, T_X) = 0$, then $X$ is \emph{rigid}: every deformation of $X$ is trivial near $0$.
\end{enumerate}
\end{theorem}

This is \cite[Chapter 6]{huybrechts-complex} and \cite[Chapter 9]{voisin-hodge-1}. For curves, $H^2(X, T_X) = 0$ for dimension reasons, so curves are unobstructed; with
Proposition~\ref{deformations:proposition-curves-deformations}, $\PP^1$ is rigid, an elliptic curve has a one-dimensional smooth Kuranishi space (the family over $\mathbb{H}$ near $\tau_0$), and a curve of genus $g \geq 2$
has a smooth $(3g - 3)$-dimensional Kuranishi space, universal because such curves have no holomorphic vector fields ($T_X$ is a line bundle of negative degree $2 - 2g$, Chapter~\ref{curves}).

\begin{theorem}[Bogomolov--Tian--Todorov]\label{deformations:theorem-btt}
A compact K\"ahler manifold $X$ with trivial canonical bundle is unobstructed: its Kuranishi space is smooth, of dimension $h^1(X, T_X) = h^{n-1,1}(X)$.
\end{theorem}

This is \cite[Chapter 6]{huybrechts-complex}. The equality $h^1(T_X) = h^{n-1,1}$ comes from the isomorphism $T_X \cong \Omega^{n-1}_X$ given by contraction with a nowhere vanishing holomorphic $n$-form. It applies
to Calabi--Yau manifolds (Chapter~\ref{calabi-yau}) and to K3 surfaces, whose Kuranishi space has dimension $h^{1,1} = 20$.

\begin{theorem}[Kodaira--Spencer]\label{deformations:theorem-kahler-stability}
Let $\pi \colon \mathcal{X} \to B$ be a family. If $X_0$ is K\"ahler, then $X_t$ is K\"ahler for $t$ in a neighbourhood of $0$. The Hodge numbers $h^{p,q}(X_t)$ are upper semicontinuous functions of $t$ for any
family, and constant near $0$ if $X_0$ is K\"ahler.
\end{theorem}

This is \cite[Chapter 9]{voisin-hodge-1}. Constancy for K\"ahler fibres follows from semicontinuity and the Hodge decomposition: $\sum_{p+q=k}h^{p,q}(X_t) = b_k$ is constant by
Corollary~\ref{deformations:corollary-fibres-diffeomorphic}, and a sum of upper semicontinuous functions with constant sum consists of locally constant functions. Being projective, in contrast, is not an open
condition: a general small deformation of a projective K3 surface is not projective.

\begin{remark}\label{deformations:remark-jumping}
\emph{Jumping of the complex structure.} The Hirzebruch surfaces $\Sigma_n = \PP(\mathcal{O} \oplus \mathcal{O}(n))$ over $\PP^1$ are diffeomorphic for all even $n$, and there is a family over $\CC$ whose fibre over $t \neq 0$ is
$\Sigma_0 = \PP^1 \times \PP^1$ and whose fibre over $0$ is $\Sigma_2$ \cite{kodaira-deformations}. So isomorphism classes of fibres can jump at special points, and the set of isomorphism classes of complex
structures on a fixed manifold is in general not a Hausdorff space: the point $[\Sigma_2]$ lies in the closure of the point $[\Sigma_0]$. Correspondingly $h^0(T_{\Sigma_2}) = 7 > 6 = h^0(T_{\Sigma_0})$: the
special fibre has more symmetry.
\end{remark}

\begin{remark}\label{deformations:remark-moduli-curves}
\emph{Moduli of curves.} The compact Riemann surfaces of genus $g \geq 2$ form a coarse moduli space $\mathcal{M}_g$, a normal quasi-projective variety of dimension $3g - 3$, as Riemann predicted; it is
the quotient of the \emph{Teichm\"uller space} $\mathcal{T}_g$, which classifies complex structures together with a marking of the fundamental group and is biholomorphic to a bounded domain in $\CC^{3g - 3}$,
by the action of the \emph{mapping class group}; its singularities come from curves with nontrivial automorphisms \cite{harris-morrison}. The \emph{period map} sends a curve to its Jacobian
(Chapter~\ref{abelian-varieties}); by Torelli's theorem it is injective, and its derivative is computed from the Kodaira--Spencer map by cup product, a basic instance of the variations of Hodge structure
of Chapter~\ref{vhs}.
\end{remark}

\section{Exercises}\label{deformations:section-exercises}

\begin{exercise}\label{deformations:exercise-hopf}
Show that every Hopf surface $X_\alpha$ (Example~\ref{deformations:example-families}(3)) is diffeomorphic to $S^1 \times S^3$.
\end{exercise}

\begin{solution}
Write $\alpha = re^{i\theta}$ with $r > 1$. The map $\CC^2 \setminus \{0\} \to \RR \times S^3$, $z \mapsto (\log|z|/\log r, z/|z|)$, is a diffeomorphism transforming $z \mapsto \alpha z$ into
$(s, w) \mapsto (s + 1, e^{i\theta}w)$. So $X_\alpha$ is the quotient of $\RR \times S^3$ by the group generated by $(s, w) \mapsto (s + 1, e^{i\theta}w)$, the \emph{mapping torus} of the rotation $w \mapsto e^{i\theta}w$ of $S^3 \subseteq \CC^2$.
The map $[s, w] \mapsto ([s], e^{-i\theta s}w)$, with $[s] \in \RR/\ZZ$, is well defined from this quotient to $S^1 \times S^3$: replacing $(s, w)$ by $(s + 1, e^{i\theta}w)$ gives
$e^{-i\theta(s + 1)}e^{i\theta}w = e^{-i\theta s}w$. In the same way $([s], u) \mapsto [s, e^{i\theta s}u]$ is well defined, and the two maps are smooth and inverse to each other.
\end{solution}

\begin{exercise}\label{deformations:exercise-automorphisms-elliptic}
Show that the group of automorphisms of $E_\tau$ fixing $0$ is $\{\pm1\}$ unless $\tau$ is equivalent to $i$ or $\rho$, where it is cyclic of order $4$ or $6$.
\end{exercise}

\begin{solution}
By Proposition~\ref{deformations:proposition-tori-isomorphic}, such an automorphism is $z \mapsto \alpha z$ with $\alpha\Lambda = \Lambda$, $\Lambda = \ZZ + \ZZ\tau$. Then $|\alpha| = 1$ (the shortest nonzero vectors of $\Lambda$ are
permuted) and $\alpha$ is a root of the characteristic polynomial of an integer matrix of size $2$, so $\alpha + \bar\alpha \in \ZZ$, and $\alpha \in \{\pm1, \pm i, \pm\rho, \pm\rho^2\}$. If $\alpha = \pm i$, then
$i\Lambda = \Lambda$, so $\Lambda$ is closed under multiplication by $\ZZ[i]$. Let $\lambda \in \Lambda$ be a nonzero element of smallest absolute value. Then $\lambda\ZZ[i] \subseteq \Lambda$; conversely, for
$\mu \in \Lambda$ there is $q \in \ZZ[i]$ with $|\mu/\lambda - q| < 1$ (round the real and imaginary parts, as in the division with remainder of Example~\ref{rings:example-gaussian-integers}), so
$\mu - q\lambda \in \Lambda$ is shorter than $\lambda$, hence $0$. So $\Lambda = \lambda\ZZ[i]$, $E_\tau \cong E_i$, and $\tau$ is equivalent to $i$ by Theorem~\ref{deformations:theorem-elliptic-moduli}. If $\alpha$ is a
primitive third or sixth root of unity, then $\rho$ is a power of $\alpha$, so $\rho\Lambda = \Lambda$, and the same argument with the Eisenstein integers $\ZZ[\rho]$, which have a division with remainder
of the same kind (Proposition~\ref{number-theory:proposition-eisenstein}, where $\rho$ is written $\omega$), gives $\tau \sim \rho$. Conversely $i\ZZ[i] = \ZZ[i]$ and $\rho\ZZ[\rho] = \ZZ[\rho]$.
\end{solution}

\begin{exercise}\label{deformations:exercise-legendre-family}
For $\lambda \in \CC \setminus \{0, 1\}$ let $C_\lambda$ be the cubic curve $y^2 = x(x - 1)(x - \lambda)$, completed in $\PP^2$. Show that the substitutions $x \mapsto 1 - x$ and $x \mapsto x/\lambda$ (with suitable changes
of $y$) give isomorphisms $C_\lambda \cong C_{1 - \lambda}$ and $C_\lambda \cong C_{1/\lambda}$, so that the six values $\lambda$, $1 - \lambda$, $1/\lambda$, $1/(1 - \lambda)$, $\lambda/(\lambda - 1)$, $(\lambda - 1)/\lambda$
give isomorphic curves, and check that
\[
J(\lambda) = \frac{(\lambda^2 - \lambda + 1)^3}{\lambda^2(\lambda - 1)^2}
\]
takes the same value at these six points. (In fact $j = 256J$ is the $j$-invariant, and $C_\lambda \cong C_\mu$ if and only if $J(\lambda) = J(\mu)$.)
\end{exercise}

\begin{solution}
With $x = 1 - x^{\prime}$ and $y = iy^{\prime}$: $-y^{\prime 2} = (1 - x^{\prime})(-x^{\prime})(1 - x^{\prime} - \lambda)$, that is, $y^{\prime 2} = x^{\prime}(x^{\prime} - 1)(x^{\prime} - (1 - \lambda))$. With $x = \lambda x^{\prime}$ and
$y = \lambda^{3/2}y^{\prime}$: $\lambda^3y^{\prime 2} = \lambda x^{\prime}(\lambda x^{\prime} - 1)(\lambda x^{\prime} - \lambda) = \lambda^3x^{\prime}(x^{\prime} - 1/\lambda)(x^{\prime} - 1)$. The two maps $\lambda \mapsto 1 - \lambda$ and
$\lambda \mapsto 1/\lambda$ generate a group of order $6$ acting on $\CC \setminus \{0, 1\}$, with the six listed values as orbit. For invariance it suffices to check the generators: $\lambda^2 - \lambda + 1$ and
$\lambda(\lambda - 1)$ are invariant under $\lambda \mapsto 1 - \lambda$; under $\lambda \mapsto 1/\lambda$, the numerator becomes $\lambda^{-6}(1 - \lambda + \lambda^2)^3$ and the denominator
$\lambda^{-2}\lambda^{-2}(1 - \lambda)^2 = \lambda^{-4}(\lambda - 1)^2$, so $J(1/\lambda) = \lambda^{-6}\lambda^4(\lambda^2 - \lambda + 1)^3/(\lambda - 1)^2 = J(\lambda)$.
\end{solution}

\begin{exercise}\label{deformations:exercise-parameter-counts}
Count parameters, as in Example~\ref{deformations:example-families}(2), for the following families of curves, and compare with Riemann's count $3g - 3$:
\begin{enumerate}
\item smooth plane quartics, of genus $3$, modulo projective transformations of $\PP^2$;
\item the curves $y^2 = f(x)$ with $f$ a polynomial with $2g + 2$ distinct roots (the \emph{hyperelliptic curves} of genus $g$), which are determined by the set of roots of $f$ on $\PP^1$ modulo
the M\"obius transformations of $\PP^1$.
\end{enumerate}
\end{exercise}

\begin{solution}
(1) Quartic forms in three variables have $\binom{6}{2} = 15$ coefficients; subtracting $1$ for scaling and $\dim\mathrm{PGL}_3(\CC) = 8$ gives $6 = 3 \cdot 3 - 3$. (2) $2g + 2$ points of $\PP^1$, modulo the group
$\mathrm{PGL}_2(\CC)$ of dimension $3$, give $2g + 2 - 3 = 2g - 1$ parameters. For $g = 2$ this is $3 = 3g - 3$, and in fact every curve of genus $2$ is hyperelliptic. For $g \geq 3$, $2g - 1 < 3g - 3$, so
the hyperelliptic curves form a proper subfamily: a general curve of genus $g \geq 3$ is not hyperelliptic.
\end{solution}
